%% ============================================================
%% Robot-control research paper — IEEEtran template
%% (ICRA / IROS / RA-L / T-RO style). Build with:
%%     ./build.sh        (bash)     or
%%     ./build.ps1       (Windows PowerShell)   or
%%     make paper        (from project root)
%%
%% Class options:
%%   [conference]      -> ICRA/IROS conference format (default below)
%%   [journal]         -> T-RO / TASE journal format
%%   For RA-L use [letterpaper, journal] and the RA-L instructions.
%% ============================================================
\documentclass[conference]{IEEEtran}

\usepackage[utf8]{inputenc}
\usepackage[T1]{fontenc}
\usepackage{amsmath,amssymb,amsfonts}
\usepackage{graphicx}
\usepackage{booktabs}
\usepackage{algorithmic}
\usepackage{url}
\usepackage[colorlinks=true,allcolors=blue]{hyperref}

% Look for figures in this folder, then in the generated results.
\graphicspath{{figures/}{../../results/figures/}{../../results/figures/run_01/}}

% --- Title & authors --------------------------------------------------
\title{<Paper Title>:\\ A Robot-Control Study}

\author{
  \IEEEauthorblockN{First Author\IEEEauthorrefmark{1}}
  \IEEEauthorblockA{\IEEEauthorrefmark{1}Department of XXX, University YYY\\
  Email: author@example.com}
}

\begin{document}
\maketitle

% --- Abstract & keywords ----------------------------------------------
\begin{abstract}
% One paragraph: problem, approach, key result, significance.
This paper presents \dots. We propose \dots. Simulation results on a
double-pendulum benchmark show \dots.
\end{abstract}

\begin{IEEEkeywords}
robot control, dynamics, computed torque, LQR, simulation
\end{IEEEkeywords}

% --- 1. Introduction --------------------------------------------------
\section{Introduction}
\label{sec:intro}
Background and motivation. State the research question and contributions:
\begin{itemize}
  \item Contribution 1.
  \item Contribution 2.
\end{itemize}

% --- 2. Related Work --------------------------------------------------
\section{Related Work}
\label{sec:related}
Prior work and how this study differs. Cite e.g.~\cite{featherstone2008rigid}.

% --- 3. Methodology ---------------------------------------------------
\section{Methodology}
\label{sec:method}
We model the robot by the rigid-body equation of motion
\begin{equation}
  M(q)\,\ddot{q} + C(q,\dot{q})\,\dot{q} + g(q) = \tau,
  \label{eq:eom}
\end{equation}
with state $x = [q^\top, \dot{q}^\top]^\top$ and joint torque input $\tau$.
The forward dynamics are computed with the Articulated Body Algorithm
(Pinocchio \texttt{aba}) and integrated with a fixed-step Runge--Kutta (RK4)
scheme. Controller design (e.g.\ computed torque) follows from~\eqref{eq:eom}.

% --- 4. Experiments ---------------------------------------------------
\section{Experiments}
\label{sec:exp}
Experimental setup: robot model, initial conditions, controller gains
(see \texttt{config/params.yaml}), and evaluation metrics. Each experiment is
documented in \texttt{docs/experiments/exp\_NN\_*.md}.

% Example figure (uncomment after running an experiment that writes the PNG):
% \begin{figure}[t]
%   \centering
%   \includegraphics[width=\columnwidth]{free_response.png}
%   \caption{Free response of the double pendulum (EXP-01).}
%   \label{fig:free_response}
% \end{figure}

% --- 5. Results & Discussion ------------------------------------------
\section{Results and Discussion}
\label{sec:results}
Report quantitative results.

\begin{table}[t]
  \centering
  \caption{Closed-loop performance metrics.}
  \label{tab:metrics}
  \begin{tabular}{lcc}
    \toprule
    Metric & Baseline & Proposed \\
    \midrule
    Settling time [s] & -- & -- \\
    RMSE              & -- & -- \\
    Control effort    & -- & -- \\
    \bottomrule
  \end{tabular}
\end{table}

% --- 6. Conclusion ----------------------------------------------------
\section{Conclusion}
\label{sec:conclusion}
Summary, limitations, and future work.

% --- Bibliography -----------------------------------------------------
% Uses the shared bib file at the repository root: references/references.bib
\bibliographystyle{IEEEtran}
\bibliography{../../references/references}

\end{document}
